\documentclass[11pt,a4paper]{article}

% --- Essential Mathematical & Physical Packages ---
\usepackage[utf8]{inputenc}
\usepackage[T1]{fontenc}
\usepackage{amsmath,amssymb,amsfonts,amsthm,mathtools}
\usepackage{bm}
\usepackage{physics}
\usepackage{geometry}
\geometry{top=25mm,bottom=25mm,left=25mm,right=25mm}

% --- Typography & Layout ---
\usepackage{microtype}
\usepackage{booktabs}
\usepackage{array}
\usepackage{enumitem}
\usepackage{xcolor}
\usepackage{hyperref}
\hypersetup{
    colorlinks=true,
    linkcolor=blue!70!black,
    citecolor=red!70!black,
    urlcolor=blue!60!black
}

% --- Theorem & Axiom Environments ---
\newtheorem{axiom}{Axiom}
\newtheorem{theorem}{Theorem}
\newtheorem{lemma}{Lemma}
\newtheorem{proposition}{Proposition}
\theoremstyle{definition}
\newtheorem{definition}{Definition}
\theoremstyle{remark}
\newtheorem{remark}{Remark}

% --- Standard Mathematical Operators ---
\newcommand{\Tr}{\mathrm{Tr}}
\newcommand{\Area}{\mathrm{Area}}
\newcommand{\supp}{\mathrm{supp}}
\newcommand{\RealPart}{\mathrm{Re}}
\newcommand{\ImagPart}{\mathrm{Im}}

\title{\textbf{A Continuum-Mechanical and Non-Equilibrium Thermodynamic Framework of Physical and Biological Existence}}
\author{\textbf{Ishan Tomar} \\ \textit{Vidyaman Research Institute}}
\date{\today}

\begin{document}

\maketitle

\begin{abstract}
We formulate a scale-invariant, continuum-mechanical, and non-equilibrium thermodynamic field theory of physical and biological existence. Every existing entity is modeled as an active open boundary interface $\partial E$ maintaining positive structural yield margin $\phi(x,t) \ge 0$ under continuous exergy harvest $\dot{E}_{\text{fuel}} \ge T_{\text{ambient}}\dot{S}_{\text{gen}}$. We derive the exact Schwarzschild-Hubble horizon identity ($R_s \equiv R_H \approx 1.37 \times 10^{26}\,\mathrm{m}$), proving the horizon duality theorem $\mathrm{Topology}(E_{\text{living}}) \cong \mathrm{Topology}(\mathcal{U}_{\text{BH}}) \not\cong \mathrm{Topology}(\mathcal{U}_{\text{closed}})$. Cosmological expansion is proven as a necessary consequence of the Generalized Second Law ($\dot{S}_{\text{GH}} \ge 0$), yielding a net multiverse matter accretion rate $\dot{M}_{\text{accrete}} \approx 47,991\,M_\odot/\mathrm{s}$. Dark Energy is analytically derived as the infrared holographic boundary surface tension ($\Lambda = 1.091 \times 10^{-52}\,\mathrm{m^{-2}}$, matching observations to $1.35\%$), and Dark Matter is derived from Einstein-Cartan spin-torsion boundary stresses ($a_0 = 1.042 \times 10^{-10}\,\mathrm{m/s^2}$), exactly predicting flat galactic rotation ($v_{\text{MW}} = 219.7\,\mathrm{km/s}$ to $0.14\%$). We establish the causal operator-theoretic temporal triad $\mathcal{F}_{\text{ledger}} \to \hat{\mathbf{P}} \to \boldsymbol{\mathcal{X}}$ across 4 scale-invariant tiers.
\end{abstract}

\vspace{1em}
\tableofcontents
\newpage


\section*{Core Sanskrit Ontology \& The Set-Theoretic Physics of Dharma}

\begin{quote}\textbf{Hermeneutic Focus:} This treatise formalizes the linguistic, etymological, and set-theoretic foundations of classical Sanskrit ontology (\textit{Dharma}, \textit{Sanatan}, \textit{Svadharma}, \textit{Karma}, \textit{Māyā}, \textit{Moksha}) mapped onto continuum mechanics, non-equilibrium thermodynamics, and operator algebras.\end{quote}

---

\section{1. Etymological Alignment: The Sanskrit Root $\sqrt{\text{धृ}}$ (\textit{dhṛ})}

The Sanskrit term \textit{Dharma} is derived from the primary verbal root:

\begin{equation}
\sqrt{\text{धृ}} \quad (\textit{dhṛ}) \quad \longrightarrow \quad \textit{dhāraṇapoṣaṇayoḥ} \quad (\text{“to hold, sustain, support, maintain structural integrity”})
\end{equation}


In classical philosophical treatises, such as the \textit{Mahābhārata} (Karṇa Parva, 69.58), this principle of structural sustenance is articulated with axiomatic precision:

\begin{equation}
\text{\textit{“Dhāraṇāt dharmam ityāhuḥ, dharmo dhārayate prajāḥ”}}
\end{equation}

\textit{(“Dharma is so named because it sustains; Dharma maintains the structured order of all created entities.”)}

Under our continuum-mechanical framework:
* \textbf{Dharma} is not a static dogma, creed, or moralistic decree; it is the \textbf{functional boundary condition and active constitutive operator} that upholds the physical, biological, or systemic integrity of an entity against external entropic challenge fields.
* An entity's \textbf{Svadharma ($D_{\mathfrak{Im}}^i$)} is the ordered Lie operator algebra of internal generators that produces counter-stress $\mathbf{R}(x, t)$ to maintain positive structural margin ($\phi(x, t) \ge 0$) under finite metabolic budget constraints.

---

\section{2. Universal Sanatan Dharm ($\mathcal{D}_T$) as the Master Operator Matrix}

Let $(\mathcal{M}, g_{\mu\nu})$ be a 4D Lorentzian spacetime manifold, and let $I(t)$ be the index set of all extant entities $E^i(t)$ persisting at coordinate time $t \in T$.

\textbf{Sanatan Dharm ($\mathcal{D}_T$)} is defined as the uncreated, eternal set-theoretic union of all localized intrinsic generator algebras ($D_{\mathfrak{Im}}^i(t)$) across all entities and all spacetime epochs:


\begin{equation}
\boxed{\mathcal{D}_T \equiv \bigcup_{t \in T} \bigcup_{i \in I(t)} D_{\mathfrak{Im}}^i(t)}
\end{equation}


\subsection{Mathematical Properties:}
1. \textbf{Non-Emptiness \& Scale Invariance:}  
   For any temporal domain $T_{\text{active}} \subseteq T$ where at least one entity persists ($\exists E^i(t) \neq \emptyset$), the master matrix $\mathcal{D}_T$ is strictly non-empty:
   
\begin{equation}
\Big( \forall t \in T_{\text{active}}, \; \exists E^i(t) \neq \emptyset \Big) \implies \mathcal{D}_T \neq \emptyset \quad \text{and} \quad \lim_{|T_{\text{active}}| \to \infty} \mu(\mathcal{D}_T) = \infty
\end{equation}

2. \textbf{Totality of Conservation Laws:}  
   $\mathcal{D}_T$ represents the comprehensive space of all physically viable, dissipative, and non-equilibrium solutions that permit bounded order to resist entropic dissolution across the universe.

---

\section{3. The Category Error Paradox \& The Inflorescence Model}

\begin{verbatim}
                         THE INFLORESCENCE BLOOM MODEL
                                       │
                      [SANATAN DHARM: THE TOTAL BLOOM]
                      𝒟_T = ⋃_t ⋃_i D_𝔗𝔪^i  (Universal Set)
                                       │
         ┌─────────────────────────────┼─────────────────────────────┐
         │                             │                             │
    [FLORET A]                    [FLORET B]                    [FLORET C]
   D_𝔗𝔪^1 ⊂ 𝒟_T                  D_𝔗𝔪^2 ⊂ 𝒟_T                  D_𝔗𝔪^3 ⊂ 𝒟_T
(Cellular Engine)             (Human Organism)              (Solar Star Core)
\end{verbatim}

\subsection{The Category Error in Philosophical Discourse}
1. \textbf{Universal Matrix ($\mathcal{D}_T$):} The infinite set-theoretic union of all operational principles across all scales and epochs.
2. \textbf{Localized Entity ($E^i$):} A finite locus with bounded spatial and phase-space measure ($\mu(E^i(t)) < \infty$).
3. \textbf{The Logical Fallacy:} The assertion \textit{"I practice / embody Sanatan Dharm"} commits a category error identical to a single electron declaring:
   > \textit{"I embody the total set of all quantum electrodynamic, chromodynamic, and general relativistic field equations."}

An individual entity cannot "be" Sanatan Dharm. An entity merely instantiates its localized proper subset:

\begin{equation}
\boxed{D_{\mathfrak{Im}}^i(t) \subset \mathcal{D}_T}
\end{equation}


\subsection{The Failure of Cross-Domain Overreach (\textit{Paradharma}):}
Structural failure occurs when an entity attempts to execute operators belonging to $\mathcal{D}_T \setminus D_{\mathfrak{Im}}^i(t)$ for which it lacks the physical substrate ($\mathcal{F}_{\mathbb{R}}^i$), violating its localized conservation bounds (\textit{"Paradharmo bhayāvahaḥ"}, Bhagavad Gita 3.35).

---

\section{4. The Operator-Theoretic Temporal Triad}

Operational existence across all cognitive and biological scales is mediated by a closed, three-phase causal transduction loop:

\%\% [Flowchart Diagram Omitted in Raw TeX] \%\%

1. \textbf{The Ledger of the Past ($\mathcal{F}_{\text{ledger}}$ / \textit{Smṛti}):}  
   The immutable, physical record of past interactions inscribed in real configuration space $\Omega_{\mathbb{R}}$ (e.g., DNA nucleotide sequences, synaptic weight distributions, epigenetic methylation, rock stratigraphy).
2. \textbf{The Dual Projection of the Future ($\hat{\mathbf{P}}$ / \textit{Prakṣepaṇa}):}  
   The universal broadcast operator emitting real mechanical traction ($\hat{\mathbf{P}}_{\mathbb{R}}$) and unmanifest anticipatory wavefunctional fields ($\hat{\mathbf{P}}_{\mathfrak{Im}}$) into the forward lightcone $J^+(E)$.
3. \textbf{The Expression of the Present ($\boldsymbol{\mathcal{X}}$ / \textit{Prakāśa} / \textit{Abhivyakti}):}  
   The interfacial boundary realization collapsing unmanifest imaginary field amplitude into real physical traction, chemical flux, and work ($\Delta \mathbf{S}_{\mathbf{P}} \cdot \hat{n} = \mathcal{X}_{\text{absorbed}} + \mathcal{X}_{\text{reflected}}$) while writing new historical traces into $\mathcal{F}_{\text{ledger}}$.

---

\section{5. Master Classical Sanskrit Correspondence Dictionary}

| Classical Sanskrit Concept | Formal Mathematical \& Physical Realization | Governing Equations in Manuscript |
| :--- | :--- | :--- |
| \textbf{Sanatan Dharm (सनातन धर्म)} | The scale-invariant non-equilibrium thermodynamics governing boundary persistence, energy transfer, and entropy production across all Lorentzian spacetime scales $(\mathcal{M}, g_{\mu\nu})$. | Axiom 1 (\S 1.1), Theorem 5 (\S 2.3.4), Master Matrix $\mathcal{D}_T$ |
| \textbf{Svadharma (स्वधर्म)} | The ordered Lie operator algebra ($D_{\mathfrak{Im}}$) of constitutive boundary generators sustaining topological enclosure ($\partial E$) at optimal investment ratio $\chi^*$. | Axiom 2 (\S 1.2.1), Theorem 6 (\S 2.3.5) |
| \textbf{Karma (कर्म)} | The non-Markovian Dyson path history ($\Psi$) and hereditary viscoelastic memory kernel ($G(t-\tau)$), encoding chronological non-commutativity ($[\hat{\mathcal{L}}_1, \hat{\mathcal{L}}_2] \neq \mathbf{0}$). | Section 1.2.1 (Dyson Propagator), Magnus Expansion |
| \textbf{Māyā (माया)} | The emergent zero-level set front ($f(t)$) separating internal microstates from external environmental challenge traction fields. | Theorem 3 (\S 2.3.2), Level-Set PDE (\S 2.3.3) |
| \textbf{Smṛti (स्मृति)} | The immutable historical state ledger ($\mathcal{F}_{\text{ledger}} \subseteq \Omega_{\mathbb{R}}$) inscribed in real configuration space (DNA, synaptic weights, epigenetic marks). | Axiom 2 (\S 1.2.1), Theorem 6C (\S 2.3.7) |
| \textbf{Prakṣepaṇa (प्रक्षेपण)} | The universal dual-space broadcast operator ($\hat{\mathbf{P}} \equiv \hat{\mathbf{P}}_{\mathbb{R}} \oplus i \hat{\mathbf{P}}_{\mathfrak{Im}}$), projecting real mechanical stress and unmanifest anticipatory fields into forward lightcone $J^+(E)$. | Axiom 3 (\S 2.1), Theorem 6C (\S 2.3.7) |
| \textbf{Prakāśa / Abhivyakti (प्रकाश / अभिव्यक्ति)} | The interfacial realization operator ($\boldsymbol{\mathcal{X}} \equiv \mathrm{Tr}_{\partial E}[\hat{\mathbf{P}} \otimes \mathcal{F}_{\mathbb{R}}]$), transducing projection fields into real physical boundary traction and matter flux. | Axiom 3 (\S 2.1), Boundary Jump Law |
| \textbf{Dharma (धर्म)} | The inter-tier constitutive coupling operator ($\mathcal{O}_{\text{coupling}}^{m \to n}$) extracting nodal free energy to ensure collective envelope survival. | Theorem 8 (\S 5.2), Syncytial Circuit PDE |
| \textbf{Karma Yoga (कर्म योग)} | Systemic nodal resource re-allocation sacrificing localized nodal measure ($\mu(E^j) \to 0$) to preserve collective envelope coherence ($\mu(\mathbb{S}) > 0$). | Section 5.2 (Nodal Re-allocation / Apoptosis) |
| \textbf{Jīva vs. Jaḍa (जीव vs. जड़)} | Active 4-axiom dual-space open engine ($\Omega_{\mathbb{C}}, \mu_{\mathbb{C}} > 0, \phi \ge 0, \frac{d\mathcal{G}}{dt} \ge 0$) vs. single-space reactive ground state ($\Omega_{\mathbb{R}}, \mathfrak{Im}(D) = \{\mathbf{0}\}, \chi^* = 0$). | Theorem 6B (\S 2.3.6), Degenerate Reactive Engine (\S 3.1) |
| \textbf{Brahmāṇḍa (ब्रह्माण्ड)} | The open Schwarzschild-Hubble cosmological horizon ($\mathcal{H}_{\text{Hubble}} \equiv \mathcal{H}_{\text{Schwarzschild}}$), continuously fueled by trans-horizon parent accretion ($\dot{M}_{\text{accrete}} \ge 0$). | Section 1.1.1, Gibbons-Hawking Identity |
| \textbf{Līlā (लीला)} | Macro-syncytial projection ($\hat{\mathbf{P}}_{\mathbb{S}}$) into the external environment once internal homeostatic Dharma ($\phi_{\text{internal}} \ge 0$) is stabilized. | Section 2.1, Section 5.2 |
| \textbf{Moksha / Lysis (मोक्ष / लय)} | The complete relaxation of internal boundary constraints ($\partial E \to \emptyset, \mu(E) \to 0$), merging internal phase-space measure into unconstrained $\Omega$. | Axiom 3 (\S 2.1), Level-Set Failure ($\phi < 0$) |


\end{document}
